% ============================================================
%  LUCRARE DE LICENȚĂ - IWB Page Builder  (FIȘIER UNIC)
%  Universitatea de Vest din Timișoara, Facultatea de Informatică
%
%  Compilare pe Overleaf:
%    1. Uploadează DOAR acest fișier (fără altele)
%    2. Settings → Compiler → pdfLaTeX
%    3. Recompile (poate fi nevoie de 2 rulări pentru cuprins)
%
%  Compilare locală:
%    pdflatex licenta_iwb_complet.tex
%    pdflatex licenta_iwb_complet.tex   (a doua oară pentru cuprins)
% ============================================================

\documentclass[12pt,a4paper]{report}

% ── Encoding & limbă ─────────────────────────────────────────────────────────
\usepackage[utf8]{inputenc}
\usepackage[T1]{fontenc}
\usepackage[romanian]{babel}

% ── Margini și spațiere ──────────────────────────────────────────────────────
\usepackage[a4paper, top=2.5cm, bottom=2.5cm, left=3cm, right=2.5cm]{geometry}
\usepackage{setspace}
\onehalfspacing

% ── Grafică ──────────────────────────────────────────────────────────────────
\usepackage{graphicx}
\usepackage{float}
\usepackage{caption}

% ── Tabele ───────────────────────────────────────────────────────────────────
\usepackage{booktabs}
\usepackage{tabularx}
\usepackage{array}

% ── Culori ───────────────────────────────────────────────────────────────────
\usepackage{xcolor}

% ── Matematică ───────────────────────────────────────────────────────────────
\usepackage{amsmath}

% ── Enumerări personalizate ───────────────────────────────────────────────────
\usepackage{enumitem}

% ── Hyperref ─────────────────────────────────────────────────────────────────
\usepackage[hyphens]{url}   % trebuie înaintea hyperref
\usepackage[colorlinks=true, linkcolor=black, citecolor=blue!60!black,
            urlcolor=blue!60!black, unicode=true]{hyperref}
\usepackage{microtype}

% ── Header/footer ────────────────────────────────────────────────────────────
\usepackage{fancyhdr}
\pagestyle{fancy}
\fancyhf{}
\fancyhead[L]{\small\leftmark}
\fancyhead[R]{\thepage}
\renewcommand{\headrulewidth}{0.4pt}
\renewcommand{\chaptermark}[1]{\markboth{\thechapter.\ #1}{}}
\renewcommand{\sectionmark}[1]{\markright{\thesection.\ #1}}
\setlength{\headheight}{14.5pt}
\emergencystretch=1.5em
\widowpenalty=10000
\clubpenalty=10000
\renewcommand{\topfraction}{0.85}
\renewcommand{\textfraction}{0.15}
\renewcommand{\floatpagefraction}{0.9}

% ─────────────────────────────────────────────────────────────────────────────
\begin{document}

% ═══════════════════════════════════════════════════════════════════════════
%  PAGINA 1 DE TITLU
% ═══════════════════════════════════════════════════════════════════════════
\thispagestyle{empty}
\begin{center}
\includegraphics[width=90pt]{FI.png}\\[0.6cm]
{\large\bfseries UNIVERSITATEA DE VEST DIN TIMIȘOARA}\\[0.25cm]
{\large\bfseries FACULTATEA DE INFORMATICĂ}\\[0.25cm]
{\large\bfseries PROGRAMUL DE STUDII DE LICENȚĂ: Informatică}
\end{center}

\vfill

\begin{center}
{\Huge\bfseries LUCRARE DE LICENȚĂ}
\end{center}

\vfill\vfill

{\large\noindent\textbf{COORDONATOR:\hfill ABSOLVENT:}\\[0.3cm]
\noindent Asist.drd. Alexandru Munteanu \hfill Șandru Nicolae-Andrei}

\vspace{1.5cm}
\begin{center}
{\large\bfseries TIMIȘOARA\\[0.2cm] 2026}
\end{center}
\newpage

% ═══════════════════════════════════════════════════════════════════════════
%  PAGINA 2 DE TITLU
% ═══════════════════════════════════════════════════════════════════════════
\thispagestyle{empty}
\begin{center}
{\large\bfseries UNIVERSITATEA DE VEST DIN TIMIȘOARA}\\[0.25cm]
{\large\bfseries FACULTATEA DE INFORMATICĂ}\\[0.25cm]
{\large\bfseries PROGRAMUL DE STUDII DE LICENȚĂ: Informatică}
\end{center}

\vfill

\begin{center}
{\LARGE\bfseries Aplicație web de tip Page Builder\\[0.5cm]
cu arhitectură client-server, \\[0.5cm]
Django REST API și frontend SPA}
\end{center}

\vfill\vfill

{\large\noindent\textbf{COORDONATOR:\hfill ABSOLVENT:}\\[0.3cm]
\noindent Asist.drd. Alexandru Munteanu \hfill Șandru Nicolae-Andrei}

\vspace{1.5cm}
\begin{center}
{\large\bfseries TIMIȘOARA\\[0.2cm] 2026}
\end{center}
\newpage

% ═══════════════════════════════════════════════════════════════════════════
%  ABSTRACT
% ═══════════════════════════════════════════════════════════════════════════
\chapter*{Abstract}
\addcontentsline{toc}{chapter}{Abstract}

\noindent
\textit{IWB} is a full-stack web application for building web pages visually, by
dragging components rather than writing markup. This thesis describes how it was
designed and built. The architecture is the usual client-server split: a Django~5.2
backend, running on Django REST Framework, exposes a JSON API, and a Single-Page
Application written in plain HTML, CSS, and vanilla JavaScript talks to it over HTTP.
There is no front-end framework anywhere in the project (a deliberate choice,
one I defend in Chapter~3).

The backend is split into seven Django apps: \texttt{core}, \texttt{users},
\texttt{landing}, \texttt{projects}, \texttt{editor}, \texttt{exporter}, and
\texttt{payments}. Each owns one concern and nothing more. A page is not a row in a
table of widgets; it is a single JSON document, validated against 38 known component
types every time it is written. Authentication uses JSON Web Tokens, with refresh
tokens that rotate and get blacklisted after use.

The part I find most useful in practice is the export engine: it walks a project's
JSON layout and emits a self-contained HTML/CSS/JavaScript bundle, zipped and ready to
drop onto any static host, with no server required to view the result. On top of that
sit automatic versioning (a snapshot before every edit), public sharing through
unique tokens, multiple pages per project, a form-submission inbox, and paid plans
wired to Stripe.

Sixty-eight integration tests exercise the API. The modular layout was chosen so the
system could grow one app at a time; real-time collaboration, sketched in the
conclusions, is the obvious next piece.

\bigskip
\noindent\textbf{Keywords:} page builder, Django, REST API, Single-Page Application,
JSON layout, JWT authentication, static site export, Stripe.

\bigskip\hrule\bigskip

\noindent\textbf{Rezumat.}
\textit{IWB} este un editor web vizual: construiești o pagină trăgând componente pe
ecran, nu scriind cod. Lucrarea de față descrie cum l-am proiectat și cum l-am scris.
Arhitectura e clasică, client-server: un backend Django~5.2 pe Django REST
Framework, care vorbește JSON, și o aplicație Single-Page în HTML, CSS și JavaScript
fără niciun framework. Comunicarea trece doar prin HTTP, deci cele două jumătăți nu
se cunosc între ele decât prin contractul API.

Backend-ul e împărțit în șapte aplicații Django, fiecare cu o singură treabă. O
pagină nu e un set de rânduri într-o bază de date, ci un document JSON validat
înainte de scriere față de 38 de tipuri de componente. Autentificarea folosește JWT
cu refresh token care se rotește la fiecare reînnoire. Motorul de export ia layoutul
JSON și scoate un pachet HTML/CSS/JS care merge pe orice hosting static, descărcat ca
arhivă ZIP. Pe lângă astea: versionare automată, partajare publică, mai multe pagini
per proiect, colectare de formulare și plăți prin Stripe.

\bigskip
\noindent\textbf{Cuvinte-cheie:} page builder, Django, REST API, Single-Page
Application, layout JSON, autentificare JWT, export static, Stripe.

\clearpage

% ═══════════════════════════════════════════════════════════════════════════
%  CUPRINS / FIGURI / TABELE
% ═══════════════════════════════════════════════════════════════════════════
\tableofcontents\clearpage

% ═══════════════════════════════════════════════════════════════════════════
%  CAPITOLUL 1 - INTRODUCERE
% ═══════════════════════════════════════════════════════════════════════════
\chapter{Introducere}

\section{Context și descrierea lucrării}

Cine vrea azi un site își lovește repede de o problemă veche: ca să-l faci de la
zero ai nevoie de HTML, CSS și JavaScript, iar majoritatea oamenilor care au nevoie
de un site nu știu și nici nu vor să învețe aceste lucruri. Un proprietar de cafenea
vrea o pagină, nu un curs de front-end. Aici intervin editoarele vizuale de tip
\textit{page builder}: muți blocuri pe ecran și pagina se construiește singură.

IWB, de la \textit{I(t) W(ill) B(uild)}, este un astfel de
editor, dar unul pe care l-am construit eu, în browser, și pe care îl descriu în
lucrarea de față. Un utilizator logat își face un proiect, îl asamblează din
componente gata făcute, îl versionează pe parcurs și, la final, îl exportă ca
fișiere statice. Acel export e partea care contează cu adevărat: paginile rezultate
nu mai au nevoie de IWB ca să trăiască. Le pui pe orice server și funcționează.

Sub capotă e o arhitectură client-server obișnuită. Backend-ul, un API REST scris cu
Django~5.2 și Django REST Framework, livrează JSON. Frontend-ul e o
\textit{Single-Page Application} (SPA) în HTML, CSS și JavaScript nativ, și subliniez
\emph{nativ}, fiindcă am evitat intenționat orice framework. Cele două vorbesc doar
prin HTTP. Practic, poți arunca frontend-ul și scrie altul mâine; backend-ul nici nu
ar observa.

\section{Scopul și obiectivele lucrării}

Țelul a fost unul singur și destul de ambițios pentru o licență: un sistem care
acoperă tot drumul unei pagini, de la contul creat până la fișierul descărcat:
autentificare, creare și editare de proiect, versionare, export static, partajare.
Nu o demonstrație pe o singură funcție, ci un lanț întreg care merge cap
la cap.

Ca să ajung acolo, am avut câteva ținte concrete. Prima a fost să înțeleg cum se
reprezintă o pagină ca date, nu ca text: am studiat principiile REST și am ajuns la
ideea de a ține layoutul ca JSON ierarhic, ceea ce a dictat apoi tot modelul de date.
Modelul trebuia să fie suficient de elastic cât să țină 38 de tipuri de componente
și totuși să nu lase să treacă structuri invalide, de aici validarea la nivelul
API. Backend-ul l-am tăiat în șapte aplicații Django, fiecare cu granița ei.
Motorul de export a fost piesa cea mai grea: transformă layoutul JSON într-un pachet
HTML/CSS/JS care stă pe picioarele lui, fără server. Frontend-ul, SPA cu rutare pe
hash și starea de autentificare ținută în memorie, fără bibliotecă externă.
Versionarea o face automat, prin snapshot înainte de fiecare modificare, cu
restaurare la orice punct. Iar la final am legat plățile de Stripe, cu email HTML de
confirmare după activarea abonamentului.

\section{Motivația alegerii temei}

Am ales tema din două motive care s-au întâlnit. Unul ține de Django: e printre cele
mai folosite backend-uri Python \cite{djangodocs}, și mi-am dorit să-l cunosc nu
superficial, ci până la capăt. Celălalt ține de problema în sine. Un editor vizual nu
e un proiect plictisitor cu un CRUD și gata, te pune în fața unor întrebări reale:
cum reprezinți o pagină ca date ierarhice, cum o validezi înainte să o scrii, cum o
transformi înapoi în cod, cum ții starea într-un SPA fără să te încurci. Sunt
probleme care apar în multe locuri, dincolo de proiectul ăsta.

Apoi, JavaScript nativ. Puteam lua React și aș fi terminat editorul mai repede,
recunosc asta fără probleme. Am refuzat tocmai pentru că voiam să văd ce e dedesubt:
cum funcționează un router, cum sincronizezi un iframe, cum gestionezi undo/redo de
mână, înainte să las un framework să facă totul pentru mine. E o decizie pe care, la
un proiect comercial, probabil nu aș repeta-o. Pentru o licență, a fost exact ce
trebuia. Și, oricum, scheletul rezultat se poate extinde mai departe într-un produs
adevărat.

\section{Structura lucrării}

Restul lucrării e împărțit în patru capitole, plus acesta.

Înainte de a-mi descrie propria soluție, în Capitolul~2 mă uit la ce există
deja: platformele comerciale (Wix, Webflow, WordPress cu Elementor, Squarespace) și
cele open-source (GrapesJS, editoarele desktop). Le compar într-un tabel cu IWB, mai
ales pe criteriul care mă interesa cel mai tare, dacă te lasă sau nu să-ți iei
codul și să pleci.

Capitolul~3 e greul lucrării. Acolo desfac IWB bucată cu bucată: arhitectura
în trei straturi, actorii și cazurile de utilizare, cerințele, tehnologiile, modelul
de date și apoi fiecare modul de backend la rând: autentificare, proiecte, editor
și validare, motorul de export, plățile Stripe, plus frontend-ul SPA. Cine vrea să
înțeleagă cum e construit sistemul, aici găsește răspunsul.

În Capitolul~4 arăt dacă funcționează: cum am testat, cum se împart cele 68
de teste pe module, patru fluxuri urmărite de la cap la coadă, și, la fel de
important, ce \emph{nu} face sistemul încă.

Capitolul~5 strânge concluziile și pune pe masă ce ar urma: canvas vizual cu
drag-and-drop adevărat, export pe mai multe pagini, colaborare în timp real, mutarea
imaginilor în cloud și SEO pentru paginile publice.

% ═══════════════════════════════════════════════════════════════════════════
%  CAPITOLUL 2 - STATE OF THE ART
% ═══════════════════════════════════════════════════════════════════════════
\chapter{State of the Art}

\section{Platforme comerciale de tip page builder}

Piața e plină. Există deja soluții mature, cu bugete de marketing pe care un proiect
de licență nici nu și le imaginează. Nu am pretenția să le concurez, mă uit la ele
ca să văd unde se opresc și ce las eu deschis acolo unde ele închid.

\subsection{Wix}

Wix e poate cel mai cunoscut nume din zonă, cu peste 200 de milioane de conturi
\cite{wixwebsite}. Editor drag-and-drop, hosting inclus, sute de template-uri. Pe
hârtie, tot ce ai nevoie. Problema apare când vrei să pleci. Layoutul e ținut în
structuri proprii, nedocumentate, pe care nu le poți scoate afară; site-ul tău
trăiește pe serverele lor și nicăieri altundeva. Adaugă peste asta modelul SaaS cu
abonamente etajate, reclame Wix pe planul gratuit și interdicția de a lega un domeniu
propriu fără să plătești. Pentru un utilizator cât de cât tehnic, lipsa exportului de
cod e dealbreaker-ul: dacă vrei pe altă platformă, o iei de la zero. Capcana pe
care am vrut ca IWB să o evite.

\subsection{Webflow}

Webflow joacă în altă ligă. Nu e pentru proprietarul de cafenea, ci pentru designeri
care știu ce e un \texttt{flexbox} și vor control fin pe CSS, animații și interacțiuni
\cite{webflowsite}. Editorul lui îți arată în timp real structura DOM și proprietățile
CSS, seamănă mai mult cu un IDE vizual decât cu un editor de blocuri. Și, spre
deosebire de Wix, te lasă să exporți HTML/CSS, deci poți migra măcar parțial (zic
parțial fiindcă JavaScript-ul dinamic nu vine cu tine). Ce nu îți dă e un API REST
public pentru conținut, adică lucrul de la care a pornit IWB. Prețul plătit
pentru toată puterea asta: o curbă de învățare abruptă și abonamente pe măsură.

\subsection{WordPress cu Elementor}

WordPress e peste tot, vreo 43\% din site-urile lumii merg pe el \cite{w3techs},
iar Elementor pune un editor drag-and-drop deasupra. Combinația e enorm de populară,
și din motive bune. Dar are un cost arhitectural pe care îl plătești zilnic: îți
trebuie PHP și MySQL pe server, WordPress plus Elementor cer update-uri de securitate
constante (WordPress e una dintre cele mai atacate ținte de pe web, nu o spun ca
exagerare), iar fiecare plugin în plus mai taie din viteză. Layoutul Elementor stă în
tabele proprii, ca meta-date ale paginii WordPress, ceea ce face versionarea și
exportul net mai greoaie decât dacă ai pur și simplu un JSON, cum am la IWB.

\subsection{Squarespace}

Squarespace e o platformă SaaS all-in-one care mizează pe design îngrijit
\cite{squarespace}, și aici nu dezamăgește, template-urile lor arată bine. Doar că
povestea cu exportul e aceeași ca la Wix: nu poți scoate codul, ești legat de
platformă. E \textit{vendor lock-in} în formă pură. Mâine îți dublează prețul sau
îngheață o funcție de care depinzi, iar singura ta opțiune e să reconstruiești
site-ul în altă parte. ZIP-ul static al IWB e răspunsul meu la fix scenariul ăsta:
arhiva e a ta și o duci unde vrei.

\section{Soluții open-source}

\subsection{GrapesJS}

GrapesJS e cel mai apropiat ca spirit de ce am făcut eu, și totuși altceva. E un
framework open-source în JavaScript pentru \emph{a construi} editoare vizuale
\cite{grapesjs}, adică îți dă piesele, nu produsul. Backend, conturi,
autentificare: nu există, le pui tu. Diferența cu IWB e tocmai asta: GrapesJS e un
toolkit din care îți faci editorul, pe când IWB e deja aplicația întreagă, cu backend,
frontend și export la pachet.

\subsection{Editoare desktop}

Mai sunt și editoarele desktop (Pinegrow, Bootstrap Studio) care scot HTML/CSS
curat, dar se opresc acolo. Server, versionare, partajare în echipă: nimic. Pentru un
om care lucrează singur, offline, sunt perfecte. Doar că în clipa în care vrei conturi,
un API sau să strângi formulare de la vizitatorii paginii, ai ieșit din ce pot ele.

\section{Abordări arhitecturale în editoare vizuale}

Întrebarea de fond la orice page builder e: cum ții pagina pe disc? Sunt câteva
răspunsuri, și fiecare are un preț. Varianta pe care am ales-o și eu, JSON
arborescent, e cea mai naturală, fiindcă un document HTML \emph{este} deja un
arbore, și un obiect JSON se mulează pe el aproape unu-la-unu. Bonus: versionarea
devine trivială (salvezi tot arborele ca snapshot) și randarea în HTML se face printr-o
parcurgere recursivă. Editoarele desktop merg uneori pe XML sau HTML direct ca sursă
de adevăr, corect, dar incomod de procesat din cod. Iar la celălalt capăt e
normalizarea completă: spargi fiecare componentă într-un rând de tabel. Câștigi
interogări SQL puternice, dar plătești scump: structura arborescentă se bate cap în
cap cu modelul relațional, iar versionarea se complică imediat. Pentru ce face IWB,
JSON-ul a fost alegerea clară.

Tendința de fond e arhitectura \textit{headless}: un backend API-first care doar
livrează date, iar deasupra pui orice client vrei \cite{mdnFetch}. IWB e pe
linia asta: backend-ul Django e un API REST curat \cite{fielding2000}, iar
frontend-ul e o aplicație separată care s-ar putea înlocui fără ca serverul să afle.

\section{Tabel comparativ al soluțiilor existente}

Tabelul~\ref{tab:comparison} pune toate astea cap la cap, alături de IWB.

\begin{table}[htbp]
  \centering
  \caption{Compararea principalelor soluții de tip page builder}
  \label{tab:comparison}
  \renewcommand{\arraystretch}{1.3}
  \begin{tabularx}{\textwidth}{>{\bfseries}l *{5}{>{\centering\arraybackslash}X}}
    \toprule
    Criteriu & Wix & Webflow & WP + Elementor & GrapesJS & \textbf{IWB} \\
    \midrule
    Tip licență      & SaaS & SaaS  & Open+SaaS & MIT   & MIT \\
    Editor vizual    & Da   & Da    & Da        & Da    & Da \\
    Backend propriu  & Da   & Da    & PHP+MySQL & Nu    & Da (Django) \\
    Export static    & Nu   & Parț. & Nu        & Da    & Da (ZIP) \\
    Versioning       & Da   & Da    & Plugin    & Nu    & Da (auto) \\
    API REST deschis & Nu   & Limit.& Nu        & N/A   & Da \\
    Open-source      & Nu   & Nu    & Parțial   & Da    & Da \\
    Auth JWT         & Nu   & Nu    & Nu        & N/A   & Da \\
    Plăți integrate  & Da   & Da    & Plugin    & Nu    & Da (Stripe) \\
    \bottomrule
  \end{tabularx}
\end{table}

\section{Concluzii}

Concluzia care iese din tabel e simplă: niciuna dintre soluțiile open-source pe care
le-am analizat nu pune la un loc backend REST, frontend SPA și export static într-un
singur sistem coerent, documentat și testat. Ori ai toolkit fără server (GrapesJS),
ori ai editor desktop fără API, ori ai SaaS închis care nu-ți dă codul. Spațiul ăsta
gol e locul în care stă IWB, un sistem complet, modular, care exportă static și
care, pe deasupra, știe să încaseze bani prin Stripe.

% ═══════════════════════════════════════════════════════════════════════════
%  CAPITOLUL 3 - DESCRIEREA APLICAȚIEI
% ═══════════════════════════════════════════════════════════════════════════
\chapter{Descrierea Aplicației}

\section{Arhitectura sistemului}

IWB stă pe trei straturi, ca în Figura~\ref{fig:architecture}. Sus, stratul de
prezentare, un SPA \cite{mdnSPA} în HTML, CSS și JavaScript nativ, fără framework
și fără pas de build (deschizi fișierul și merge). La mijloc, creierul: un API REST
\cite{fielding2000} pe Django~5.2 care duce toată logica: validare, versionare,
autentificare, business. Jos, datele: SQLite în development, fiindcă e zero
configurație, și PostgreSQL în producție. Pe amândouă le ating doar prin ORM-ul
Django; SQL scris de mână nu există nicăieri în proiect, și asta e intenționat:
am vrut ca trecerea SQLite$\rightarrow$PostgreSQL să nu ceară nicio modificare de cod.

\begin{figure}[htbp]
  \centering
  \includegraphics[width=0.92\linewidth]{diag_architecture}
  \caption{Diagrama de arhitectură a sistemului IWB}
  \label{fig:architecture}
\end{figure}

Granița dintre prezentare și aplicație e doar HTTP: nimic altceva nu trece între ele.
De aici și libertatea de a rescrie un strat fără să-l atingi pe celălalt.
Tabelul~\ref{tab:urlmap} arată cum se împart URL-urile API pe module.

\begin{table}[htbp]
  \centering
  \caption{Harta URL-urilor API}
  \label{tab:urlmap}

  \begin{tabularx}{\textwidth}{>{\ttfamily}l X}
    \toprule
    \normalfont\textbf{Prefix URL} & \textbf{Aplicație / Responsabilitate} \\
    \midrule
    /api/             & \texttt{landing}, informații publice despre API \\
    /api/auth/*       & \texttt{users}, înregistrare, login, logout, profil \\
    /api/projects/*   & \texttt{projects}, CRUD, versioning, sharing, pagini \\
    /api/editor/projects/\{id\}/* & \texttt{editor}, operații granulare pe componente \\
    /api/export/\{id\}/ & \texttt{exporter}, export ZIP sau JSON \\
    /api/payments/*   & \texttt{payments}, Stripe checkout, status abonament \\
    \midrule
    \textit{orice altceva} & SPA catch-all, servește \texttt{index.html} \\
    \bottomrule
  \end{tabularx}
\end{table}

\section{Actori și cazuri de utilizare}

Sunt trei feluri de actori în sistem. Vizitatorul neautentificat vede pagina de
prezentare, se poate înregistra sau loga și poate deschide paginile publice ale
altora printr-un link de partajare. Atât, nimic ce ține de cont. Utilizatorul logat
e cel pentru care a fost gândit totul: creează, editează, versionează, exportă și
partajează propriile proiecte, și tot el poate cumpăra un plan Pro sau Enterprise prin
Stripe. Iar deasupra tuturor stă administratorul, care umblă la utilizatori, proiecte,
template-uri și plăți direct din Django Admin. Nu am construit o interfață separată
de administrare, fiindcă cea din Django acoperă deja tot ce aveam nevoie.

\begin{figure}[htbp]
  \centering
  \includegraphics[width=0.92\linewidth]{diag_usecase}
  \caption{Diagrama UML Use Case a aplicației IWB}
  \label{fig:usecase}
\end{figure}

Fluxurile urmăresc viața unui proiect de la naștere până la publicare. Totul începe
cu un cont: te înregistrezi cu email și parolă, te loghezi, sistemul îți dă token-uri
JWT. De aici încolo poți crea proiecte, să le listezi, să le schimbi metadatele, să
duplici unul sau să-l ștergi. Editarea propriu-zisă înseamnă adăugat, modificat, șters
și reordonat componente, și fiecare astfel de gest lasă în urmă un snapshot, fără
să ceri tu asta. Versionarea îți arată apoi tot istoricul și te lasă să sari înapoi la
orice stare. Exportul scoate fie o arhivă ZIP descărcabilă, fie un JSON de
previzualizare. Partajarea publică deschide proiectul pentru oricine are token-ul
unic. Pe proiectele cu mai multe pagini poți adăuga, reordona și alege pagina
principală. Vizitatorii paginilor publice pot completa formulare, pe care proprietarul
le citește mai târziu. Și, în fine, cumperi un plan Pro sau Enterprise printr-un
checkout Stripe care îți trimite confirmarea pe email.

\section{Cerințe funcționale și nefuncționale}

Pe partea funcțională, lucrurile se grupează firesc pe câteva zone. Autentificarea:
cont cu email unic și parolă, iar la login token-uri JWT care se rotesc la fiecare
reînnoire. Proiectele: utilizatorul logat le creează, listează, actualizează și
șterge pe ale lui, cu o regulă pe care nu o calc nicăieri: orice scriere validează
mai întâi JSON-ul layoutului și salvează un snapshot înainte, ca să poată reveni
oricând. Export și partajare: descarci ZIP sau vezi JSON-ul, respectiv aprinzi/stingi
accesul public și, dacă vrei, regenerezi token-ul (caz în care vechiul link moare pe
loc). Formularele de pe paginile publice le poate trimite oricine, dar cu un plafon de
zece submisii pe oră per IP, altfel paginile s-ar umple de spam într-o zi. Plățile
trec prin Stripe, cu un email HTML de confirmare la fiecare activare.

Cerințele nefuncționale m-au preocupat trei: securitate, modularitate,
portabilitate. Tot ce modifică date e în spatele JWT \cite{rfc7519}, iar accesul la
resursele cuiva e zăvorât de permisiunea \texttt{IsOwner}: nu vezi proiectele
altuia, punct. Fiecare aplicație Django face un singur lucru, iar testele ating
endpoint-urile principale din fiecare modul. Pe extensibilitate mi-am impus o regulă
pe care o verific în Capitolul~3.6: un tip nou de componentă nu trebuie să ceară
atingeri în mai mult de trei fișiere. Iar exportul static merge în orice browser
modern fără nicio dependență externă, ca și baza de date care funcționează la fel
pe SQLite și pe PostgreSQL.

\section{Tehnologii utilizate}

Backend-ul e Python~3.10+ pe Django~5.2 LTS. Am ales special varianta LTS,
fiindcă are suport până în 2028 \cite{djangodocs} și nu voiam să-mi expire framework-ul
la un an după ce predau lucrarea. Peste Django stă Django REST Framework~3.17
\cite{drfdocs}, de unde iau patru lucruri pe care altfel le-aș fi scris singur:
serializatorii care validează intrarea și scot modelele în JSON, ViewSet-urile și
APIView-urile pentru cereri, router-ele care generează URL-urile, și sistemul de
permisiuni pe care l-am extins cu propriile clase. JWT-ul vine din
\texttt{simplejwt}~5.5 \cite{simplejwt}, pe care l-am pus pe rotație automată plus
blacklisting. Restul sunt piese mici dar necesare: \texttt{django-cors-headers}~4.9
pentru CORS, \texttt{python-decouple}~3.8 ca să țin secretele afară din cod,
\texttt{Pillow}~11.2 pentru avatare. Stripe intră prin SDK-ul Python oficial
\cite{stripedocs}. Baza de date: SQLite \cite{sqlitedocs} la mine pe laptop,
PostgreSQL \cite{postgresqldocs} pentru orice deployment serios.

Frontend-ul e doar HTML5, CSS3 și JavaScript ES6+. Fără framework, fără bundler,
fără pas de build. Am ținut la asta de la prima linie. Singura concesie e Tailwind,
adus prin CDN, dar \emph{numai} în interfața editorului; în fișierele exportate nu
există urmă de Tailwind, fiindcă nu vreau ca site-urile generate să depindă de un CDN
extern. Comunicarea cu backend-ul o fac prin Fetch API \cite{mdnFetch}, cel nativ din
browser.

\section{Modelul de date}

Decizia de la baza întregului model a fost să țin layoutul ca un singur document JSON,
nu spart pe tabele. Nu a fost evidentă de la început: normalizarea, cu o componentă
pe rând, mi-ar fi dat interogări SQL mai bogate. Am ales totuși JSON-ul din trei
motive care, pentru un page builder, cântăresc mai greu. Întâi, arborele paginii se
mulează pe un obiect JSON \cite{rfc8259} aproape fără efort. Apoi, versionarea: un
snapshot înseamnă să copiez un câmp, nu zeci de rânduri din cinci tabele. Și, în fine,
operațiile zilnice ale unui editor nu cer join-uri complicate între componente: mai
mereu vrei tot arborele sau o subramură, nu „toate butoanele albastre din proiecte
create marțea". Validarea structurii o fac în aplicație, la frontiera API, înainte de
orice scriere, niciodată în baza de date.

\begin{figure}[htbp]
  \centering
  \includegraphics[width=0.92\linewidth]{diag_datamodel}
  \caption{Diagrama entitate-relație simplificată a modelelor IWB}
  \label{fig:datamodel}
\end{figure}

Cinci dintre modele, \texttt{Project}, \texttt{ProjectVersion}, \texttt{Page},
\texttt{Tag}, \texttt{FormSubmission} și \texttt{ProjectTemplate}, pleacă toate de
la \texttt{core.BaseModel}. Acolo am pus ce trebuie să aibă orice entitate: cheie
primară UUID generată automat (aleg UUID tocmai ca nimeni să nu poată ghici ID-ul
următor și să enumere proiectele altora) plus timestamp-urile de creare și modificare.

\texttt{users.User} extinde \texttt{AbstractUser} din Django. I-am mutat
autentificarea pe \texttt{email} (unic), i-am adăugat o biografie opțională, un avatar
care ajunge în folderul \texttt{avatars/}, câmpul \texttt{plan} (free/pro/enterprise,
implicit \texttt{free}) și \texttt{plan\_expires\_at}. Trucul util e proprietatea
\texttt{is\_pro}: returnează \texttt{True} dacă planul e Pro sau Enterprise și încă nu
a expirat. O citesc direct din obiectul user, fără query separat.

\texttt{Project} e inima modelului. Ține layoutul într-un JSONField și, lângă el, un
al doilea JSONField numit \texttt{meta}: titlu, descriere, limbă, CSS și JS
personalizate. Partajarea publică o controlez cu o pereche: \texttt{share\_token}
(UUID unic) și booleanul \texttt{is\_public}. Un proiect e \texttt{draft},
\texttt{published} sau \texttt{archived}, iar slug-ul e unic per utilizator, nu
global, fiindcă doi oameni au tot dreptul să aibă fiecare un proiect „portofoliu".
Pentru organizare, proiectele se leagă Many-to-Many de \texttt{Tag}.

Versionarea trăiește în \texttt{ProjectVersion}, și aici am ținut neapărat ca
snapshot-urile să fie imutabile: odată scrise, nu se mai ating. Fiecare versiune
are un număr care crește monoton, câmpul \texttt{layout\_snapshot}, o notă opțională
și o referință la cine a făcut modificarea (\texttt{SET\_NULL} la ștergere, ca să nu
pierd versiunea dacă dispare contul). \texttt{Page} reprezintă paginile individuale,
fiecare cu propriul \texttt{layout}, o ordine și steagul \texttt{is\_homepage}.
\texttt{Tag} e simplu: nume și o culoare hex. Iar \texttt{FormSubmission} strânge
ce trimit vizitatorii: ID-ul formularului, datele ca JSON, IP-ul și user agent-ul.

Schema layoutului are un singur câmp rădăcină. La vârf, \texttt{layout} e un obiect cu
o cheie, \texttt{"components"}, și aceasta e o listă de componente. O componentă are
două câmpuri obligatorii: \texttt{"id"} (un UUID unic în tot arborele) și
\texttt{"type"} (unul dintre cele 38 de tipuri), plus oricâte câmpuri opționale
specifice tipului: \texttt{"text"}, \texttt{"src"}, \texttt{"href"}, \texttt{"style"}
și așa mai departe. Componentele-container, \texttt{section}, \texttt{container},
\texttt{columns}, \texttt{footer}, \texttt{hero} și celelalte din mulțimea
\texttt{CHILDREN\_TYPES}, mai au și câmpul \texttt{"children"}, lista copiilor, ceea
ce le lasă să se imbrice oricât de adânc. Unicitatea ID-urilor o verific global cu un
helper, \texttt{\_collect\_ids()}, care coboară recursiv prin \texttt{children},
\texttt{items} și \texttt{columns} și strânge toți identificatorii înainte de scriere.
Detaliul important e că verific unicitatea \emph{între} niveluri, nu doar pe orizontală: un copil nu poate fura ID-ul părintelui.

\texttt{meta} are propria lui schemă. \texttt{"title"} ajunge titlul paginii
exportate, \texttt{"description"} devine meta description pentru motoarele de căutare,
\texttt{"lang"} e codul limbii (implicit \texttt{"ro"}, fiindcă majoritatea celor care
ar folosi IWB scriu în română), iar \texttt{"custom\_css"} și \texttt{"custom\_js"}
sunt injectate la coada fișierelor generate, pentru cine vrea să intervină manual.
Restul sunt câmpuri de temă: \texttt{"font"}, \texttt{"primaryColor"},
\texttt{"textColor"}, \texttt{"bgColor"}, \texttt{"maxWidth"}. Fiecare se duce
direct într-o variabilă CSS globală a exportului. Schimbi o culoare în \texttt{meta} și
tot site-ul se schimbă, fără să atingi nicio componentă în parte.

\section{Backend-ul Django}

\subsection{Aplicațiile \texttt{landing} și \texttt{core}}

\texttt{landing} e cât se poate de mic: un singur endpoint public, \texttt{GET /api/},
care întoarce un index JSON cu toate rutele API-ului. E documentație vie și totodată
ușa de intrare pentru un client nou: intri pe \texttt{/api/} și vezi de unde să
începi, fără să fii logat și fără să citești cod.

\texttt{core} nu are URL-uri deloc; e o bibliotecă internă pe care o folosesc toate
celelalte. Aici stau \texttt{BaseModel} (UUID + timestamps), permisiunea
\texttt{IsOwner} și, cel mai important, \texttt{validators.py}, cu
\texttt{validate\_layout()} și \texttt{validate\_component()}. Prima verifică trei
lucruri: că layoutul are cheia \texttt{"components"} cu o listă, că toate ID-urile din
arbore sunt unice (le strânge recursiv din \texttt{children}, \texttt{items} și
\texttt{columns}) și că fiecare componentă trece de \texttt{validate\_component()}. A
doua se uită la o componentă singură: are \texttt{"id"} și \texttt{"type"}? E tipul
printre cele 38 cunoscute? Și apoi coboară în copii și repetă. E codul pe care îl
consider cel mai critic din tot backend-ul, fiindcă el e zidul dintre date curate și
date corupte. De aceea are și cele mai multe teste, unsprezece, cât toată aplicația
\texttt{users}.

\subsection{Autentificare și gestionarea utilizatorilor}

\texttt{users} duce tot ciclul de viață al unui cont, și am avut grijă peste tot să
economisesc apeluri inutile. \texttt{RegisterView} creează utilizatorul și îi dă
token-urile JWT pe loc. Nu te pune să faci login imediat după ce te-ai înregistrat,
fiindcă e o frecare fără rost. \texttt{LoginView} merge cu un serializator custom care
bagă în răspuns și datele utilizatorului, nu doar token-urile, ca frontend-ul să-și
poată construi toată starea dintr-un singur request. \texttt{logout\_view}
blacklistează refresh token-ul și încheie sesiunea. Iar \texttt{ProfileView} și
\texttt{ChangePasswordView} se ocupă de profil, respectiv de schimbarea parolei (cu
verificarea parolei curente, evident, altfel cine îți fură sesiunea ți-ar putea
schimba parola direct).

JWT-ul \cite{rfc7519} l-am configurat cu rotație și blacklisting. La fiecare reînnoire
de access token, refresh-ul vechi ajunge pe o listă neagră din baza de date, așa că
nu mai poate fi folosit a doua oară. Chiar dacă cineva l-a interceptat, fereastra
lui de viață e deja închisă. Duratele celor două token-uri le țin în variabile de
mediu, ca să le pot strânge sau lărgi fără să recompilez nimic.

Concret, blocul \texttt{SIMPLE\_JWT} din \texttt{settings.py} pune
\texttt{ACCESS\_TOKEN\_LIFETIME} pe 60 de minute (suprascris de
\path{ACCESS_TOKEN_LIFETIME_MINUTES}) și \texttt{REFRESH\_TOKEN\_LIFETIME} pe 7 zile
(\path{REFRESH_TOKEN_LIFETIME_DAYS}). Ca să meargă blacklisting-ul trebuie să ai
aplicația \path{rest_framework_simplejwt.token_blacklist} în \texttt{INSTALLED\_APPS}
și două tabele în baza de date: \path{token_blacklist_outstandingtoken} și
\path{token_blacklist_blacklistedtoken}. Are însă un revers pe care l-am observat
repede: tabelele astea cresc la nesfârșit dacă nu le cureți. De aici comanda
\texttt{flushexpiredtokens}, rulată periodic, care șterge token-urile deja expirate:
fără ea, baza de date s-ar îngrășa în gol cu fiecare login.

\subsection{Gestionarea proiectelor și versionare}

\texttt{ProjectViewSet} e un \texttt{ModelViewSet} DRF (adică CRUD-ul standard vine
gratis), peste care am pus câteva acțiuni proprii. \texttt{toggle-public} aprinde sau
stinge accesul fără cont. \texttt{regenerate-token} generează un \texttt{share\_token}
nou și, prin asta, omoară instant orice link vechi pe care l-ai trimis cuiva. Util
când îți dai seama că ai partajat din greșeală. \texttt{versions} listează
snapshot-urile, iar \texttt{restore-version} te duce înapoi la oricare dintre ele. Și
aici e un detaliu la care țin: restaurarea nu aruncă starea curentă, ci o salvează
mai întâi ca un snapshot nou (cu o notă pusă automat) și abia apoi suprascrie
\texttt{layout}. Practic, nu poți pierde nimic ireversibil, nici măcar dacă restaurezi
„greșit".

Tot aici stau și restul: \texttt{SharedProjectView} pentru accesul public (întoarce
layoutul după \texttt{share\_token}, dar numai dacă \texttt{is\_public=True}),
\texttt{PageDetailView} și \texttt{PageSetHomepageView} pentru paginile individuale,
\texttt{ProjectTemplateViewSet} care listează și instanțiază template-urile gata
făcute, și \texttt{FormSubmitView}, singurul endpoint deschis spre lume, deci și
singurul cu rate limiting (maximum zece submisii pe oră per IP, prin cache-ul Django).

Când creezi o pagină nouă, slug-ul se generează singur din titlu: curăț caracterele
speciale și, dacă slug-ul există deja, adaug un sufix numeric (\texttt{-1},
\texttt{-2}, și tot așa). \texttt{reorder-pages} are o gardă pe care am pus-o
intenționat strictă: lista de ID-uri primită trebuie să conțină \emph{exact} paginile
proiectului, nici una în plus, nici una lipsă, altfel refuz reordonarea. Nu vreau
să pierd sau să dublez o pagină dintr-un request prost format. \texttt{set-homepage}
stinge \texttt{is\_homepage} pe toate paginile înainte să-l aprindă pe cea aleasă, ca
să existe mereu fix una principală, niciodată zero sau două.

\texttt{duplicate} face o copie adâncă: \texttt{layout}, \texttt{meta}, statusul,
relațiile \texttt{Tag}, totul, cu un slug nou obținut prin sufixul \texttt{-copy}
(și rezolvarea coliziunilor, ca mai sus). Paginile le copiez una câte una, fiecare cu
\texttt{layout}-ul ei. Un singur lucru \emph{nu} se moștenește: partajarea. Copia
pornește cu \texttt{is\_public=False} și fără \texttt{share\_token}. Ar fi o
surpriză neplăcută ca, duplicând un proiect, să expui din greșeală și copia spre
exterior.

\subsection{Editarea vizuală și validarea layoutului}

\texttt{editor} nu are modele proprii: lucrează direct pe câmpul \texttt{layout} al
lui \texttt{Project}. Asta a fost o decizie conștientă: tot ce ține de o pagină e deja
în acel JSON, deci de ce să mai inventez tabele? Fiecare view din modul cheamă, înainte
de orice modificare, același helper de snapshot, care creează un \texttt{ProjectVersion}
cu starea de dinainte și incrementează contorul de versiune. Figura~\ref{fig:editflow}
arată fluxul întreg al unei modificări.

\begin{figure}[!htb]
  \centering
  \includegraphics[width=0.50\linewidth]{diag_editflow}
  \caption{Fluxul de editare și export al unui proiect IWB}
  \label{fig:editflow}
\end{figure}

Orice modificare de layout trece prin aceiași patru pași, în aceeași ordine. Întâi
serializatorul DRF se uită la tipuri și la câmpurile obligatorii. Apoi
\texttt{validate\_layout()} coboară recursiv prin structura JSON. Abia după ce a trecut
de amândouă, helper-ul de snapshot îngheață layoutul curent într-un
\texttt{ProjectVersion}. Și de-abia atunci scriu noua valoare. Ordinea contează:
validez înainte să salvez snapshot-ul, ca să nu umplu istoricul cu versiuni provenite
din cereri care oricum aveau să fie respinse. Răspunsul întoarce layoutul întreg,
actualizat, ca frontend-ul să-și sincronizeze starea fără un al doilea apel.

View-urile sunt cinci. \texttt{ProjectLayoutView} citește și înlocuiește layoutul cu
totul. \texttt{ComponentUpsertView} adaugă sau actualizează o componentă după ID
(de unde și „upsert", nu mă interesează dacă există deja, o pun la locul ei).
\texttt{ComponentDeleteView} șterge una. \texttt{ComponentReorderView} reordonează
după o listă de ID-uri. Iar \texttt{ComponentDefaultsView} întoarce proprietățile
implicite ale fiecărui tip. E cel pe care frontend-ul îl cheamă în clipa în care
tragi o componentă nouă pe pagină.

Tabelul~\ref{tab:components} adună cele 38 de tipuri pe cinci categorii. Aici e și
locul unde mi-am ținut promisiunea de extensibilitate din secțiunea de cerințe: ca să
adaugi un tip nou atingi exact trei fișiere, nici mai mult. Pui tipul în setul
\texttt{VALID\_COMPONENT\_TYPES} din \texttt{core/validators.py}, scrii metoda
\texttt{\_render\_<type>()} în \texttt{exporter/engine.py} și adaugi proprietățile
implicite în dicționarul \texttt{COMPONENT\_DEFAULTS} din
\texttt{editor/serializers.py}. Trei fișiere, trei modificări. Am verificat regula
de fiecare dată când am mai adăugat un tip, și a ținut.

\begin{table}[htbp]
  \centering
  \caption{Categoriile și tipurile de componente IWB}
  \label{tab:components}
  \renewcommand{\arraystretch}{1.2}
  \begin{tabularx}{\textwidth}{>{\bfseries}l X}
    \toprule
    Categorie & Tipuri \\
    \midrule
    Layout (8)      & \texttt{navbar, section, container, columns, footer,
                      divider, spacer, banner} \\
    Conținut (7)    & \texttt{text, heading, richtext, blockquote, code\_block,
                      icon, badge} \\
    Media (5)       & \texttt{image, video, embed, gallery, logo\_strip} \\
    Interactive (5) & \texttt{button, link, social\_links, form, countdown} \\
    Secțiuni (13)   & \texttt{hero, features, pricing, testimonials, faq, cta,
                      team, stats, card, cards\_grid, timeline, tabs, contact} \\
    \midrule
    \textbf{Total (38)} & \\
    \bottomrule
  \end{tabularx}
\end{table}

\subsection{Motorul de export static}

Aici e, după mine, partea cea mai interesantă a lucrării: locul unde datele redevin
cod. \texttt{exporter/engine.py} expune trei funcții: \texttt{build\_html()},
\texttt{build\_css()} și \texttt{build\_js()}. În interior, o clasă \texttt{Renderer}
parcurge recursiv arborele de componente și trimite fiecare componentă spre metoda ei
de randare, \texttt{\_render\_<type>()}. E reversul validării: una verifică
arborele, cealaltă îl transformă înapoi în HTML, amândouă coborând prin aceeași
structură.

CSS-ul pleacă de la un set de variabile globale puse în \texttt{:root} și completate
din \texttt{meta}: \texttt{--font} (familia de fonturi, cu preconectare la Google
Fonts în \texttt{<head>}, ca să nu aștepte pagina după font), \texttt{--primary}
(accentul, pe butoane și elemente interactive), \texttt{--text}, \texttt{--bg},
\texttt{--max-w} (lățimea maximă a containerelor) și \texttt{--section-py} (padding-ul
vertical al secțiunilor). Toată ideea e că schimbi un singur set de variabile și se
schimbă tot site-ul. Nu umblu prin stilurile fiecărei componente. Pe lângă CSS,
exportul scoate și meta taguri Open Graph (\texttt{og:title}, \texttt{og:description},
\texttt{og:image}), tot din \texttt{meta}, ca pagina să arate decent când o lipești
într-un mesaj sau pe o rețea socială.

Responsivitatea o rezolv cu trei praguri \texttt{@media}. La 1024 px intervin pentru
tablete: grid-urile de patru coloane se strâng la două. La 768 px e rândul
telefoanelor: două coloane devin una, iar navbar-ul ascunde linkurile și scoate
butonul hamburger. La 480 px micșorez textul prin \texttt{clamp()} și tai din
padding-ul secțiunilor, ca pe ecran mic să nu se irosească spațiu. Hamburger-ul însuși
e cât se poate de simplu: un toggle al clasei \texttt{.open} pe containerul de
linkuri, iar CSS-ul face restul cu o tranziție pe \texttt{max-height}. Fără bibliotecă,
fără JavaScript complicat.

JavaScript-ul generat are cinci comportamente, toate scrise de mână, fără bibliotecă.
Accordion-ul de FAQ deschide și închide răspunsul prin clasa \texttt{.open} pe
\texttt{.pb-faq-item}, animat tot pe \texttt{max-height}. Tab-urile leagă butoanele de
panouri prin atributul \texttt{data-tab}: apeși un buton, primește \texttt{.active}
și se arată panoul lui. Reveal-ul la scroll merge pe \texttt{IntersectionObserver},
care animă opacitatea și translația componentelor când intră în viewport, la un prag
de 15\%; am ales \texttt{IntersectionObserver} tocmai ca să nu mai ascult evenimentul
de scroll la fiecare pixel, care e risipă de procesor. Countdown-ul citește
\texttt{data-target} ca timestamp Unix și ticăie în timp real, pornit de același
observer. N-are sens să numere înainte să-l vadă cineva. Iar formularele: intercept
\texttt{submit}, strâng câmpurile cu \texttt{FormData}, le serializez ca JSON și le
trimit cu \texttt{fetch()} la endpoint-ul public de submisii, apoi afișez în pagină
succes sau eroare, fără reîncărcare.

La export ZIP, serverul împachetează arhiva în memorie (\texttt{index.html},
\texttt{styles.css}, \texttt{main.js}) și o întoarce cu tipul
\texttt{application/zip}. Un amănunt care m-a costat ceva timp de depanare: parametrul
din URL e \texttt{?type=zip}, nu \texttt{?format=zip}. DRF interceptează singur
\texttt{?format=} pentru content negotiation, așa că, dacă foloseam acel nume, cererea
mea de ZIP era „mâncată" înainte să ajungă la cod. Am redenumit parametrul în
\texttt{type} și problema a dispărut.

\subsection{Sistemul de plăți cu Stripe}

\texttt{payments} face plăți reale prin Stripe \cite{stripedocs}, pe modelul
\textit{PaymentIntents}. Fluxul are patru pași. Întâi, frontend-ul cere cu
\texttt{GET /api/payments/config/} cheia publicabilă Stripe, ca să inițializeze
Stripe.js. Apoi un \texttt{POST} la \texttt{/api/payments/create-intent/} creează pe
Stripe un \texttt{PaymentIntent} cu suma planului ales și cu metadata aferentă (plan,
email, user ID), iar răspunsul aduce \texttt{client\_secret}-ul cu care clientul va
confirma plata.

Confirmarea pe card o face Stripe.js direct în browser, și aici e poanta de
securitate: numărul cardului nu trece niciodată prin serverul meu, deci nici nu am ce
să scurg. Partea de server vine după, printr-un \texttt{POST} la
\texttt{/api/payments/confirm/}: întreb Stripe dacă \texttt{PaymentIntent}-ul chiar e
pe \texttt{succeeded}, trec printr-o gardă de idempotență pe
\texttt{stripe\_payment\_intent\_id} (unic) și, dacă totul e curat, activez planul
pentru 30 de zile și trimit emailul HTML de confirmare. Starea curentă și istoricul se
văd la \texttt{/api/payments/status/} și \texttt{/api/payments/history/}.

Emailul de confirmare l-am scris cu CSS inline, fiindcă clienții de email, mai ales
Outlook, aruncă jumătate din regulile dintr-un \texttt{<style>} extern. Are un antet
cu gradient (de la \texttt{\#1a53d6} la \texttt{\#4F7CFF}), titlul „Plată confirmată"
și un salut cu numele utilizatorului. Apoi un tabel cu rânduri alternante înșiră
detaliile: planul activat, suma (\$9.99 pentru Pro, \$29.99 pentru Enterprise),
referința de plată (primele opt caractere ale UUID-ului, cu majuscule), cardul mascat
(\texttt{**** **** **** XXXX}), data plății și data expirării. Aceasta din urmă o
pun pe verde, ca să sară în ochi cât e valabil planul. Jos, un buton spre dashboard.
Trimit și o versiune text simplă, cu aceleași date, pentru cine nu randează HTML.

Un detaliu de robustețe la care am ținut: \texttt{send\_mail()} e apelat cu
\texttt{fail\_silently=True}. Logica e că plata e mult mai importantă decât emailul.
Dacă SMTP-ul cade (și în development pică des, fiindcă rar ai un server de mail
configurat corect), nu vreau ca eșecul trimiterii să arunce o excepție tocmai după
ce am încasat banii și am activat planul în baza de date. Activarea s-a întâmplat deja;
emailul e doar curtoazie.

\texttt{Payment} ține istoricul: \texttt{user} (FK, \texttt{CASCADE} la ștergere),
\texttt{plan} (pro/enterprise), \texttt{amount} (și aici o notă: suma e în cenți, nu
în dolari, ca să nu mă lovesc niciodată de erorile de rotunjire ale numerelor cu
virgulă), \texttt{status} (completed/failed), \texttt{reference} (UUID unic generat
la creare), \texttt{card\_last4}, \path{stripe_payment_intent_id} (ID-ul Stripe, câmpul
unic pe care se sprijină idempotența) și \texttt{created\_at}. Modelul ăsta, împreună
cu \texttt{plan} și \texttt{plan\_expires\_at} de pe \texttt{User}, îmi dă verificarea
de acces Pro pe gratis: \texttt{is\_pro} răspunde direct din obiectul user încărcat
deja din token, fără să mai lovesc baza de date încă o dată.

\section{Frontend-ul SPA}

Unde să ții token-urile în frontend e una dintre acele întrebări fără răspuns perfect,
doar compromisuri. Eu am ales așa: access token-ul stă \emph{numai} în memorie,
în variabila \texttt{\_access} a singletonului \texttt{Auth}, și dispare la
reîncărcarea paginii. Refresh token-ul, în schimb, e persistent în
\texttt{localStorage}, ca sesiunea să supraviețuiască un refresh de browser, până la 7
zile. La fiecare încărcare de pagină, \texttt{tryRefresh()} încearcă să schimbe
refresh-ul pe un access nou; dacă nu reușește, șterge tot și te trimite la login.

De ce această împărțeală, și nu să pun amândouă în \texttt{localStorage}, cum e mai
simplu? Pentru că împart riscul. Un atac XSS care reușește să injecteze script poate
citi refresh token-ul din \texttt{localStorage}, asta o recunosc, nu e o apărare
perfectă, dar nu poate atinge access token-ul, care nu există nicăieri pe disc, doar
în memoria paginii curente. Cu alte cuvinte, restrâng fereastra în care un token furat
chiar e util. E un compromis, nu o garanție; dar e mai bun decât a lăsa amândouă
token-urile la îndemână.

Frontend-ul e un SPA cu rutare pe hash, fără build tools și fără framework. Inima
comunicării e \texttt{apiFetch()}, un wrapper peste \texttt{fetch()} care pune singur
header-ul \texttt{Authorization: Bearer} pe fiecare cerere autentificată. Partea care
mi-a plăcut cel mai mult să o scriu e ce face la \texttt{401}: cheamă transparent
endpoint-ul de refresh, ia un access token nou și \emph{reia} cererea originală, fără
ca utilizatorul să observe ceva. Token-ul a expirat la mijlocul unei acțiuni? Nu se
vede nimic; acțiunea pur și simplu reușește.

Router-ul ascultă \texttt{hashchange} și leagă tipare de cale (\texttt{/},
\texttt{/auth}, \texttt{/dashboard}, \texttt{/editor/:id}, \texttt{/share/:token}) de
funcții de randare. Am mers pe hash, nu pe History API, dintr-un motiv foarte practic:
History API ar fi cerut configurație de rewrite pe server pentru fiecare rută, pe când
hash-ul merge oriunde pui fișierele, fără nimic pe partea de server. Pentru un proiect
care se laudă cu deployment simplu, era alegerea coerentă.

Editorul se învârte în jurul unui panou lateral cu componentele grupate pe categorii.
Când adaugi una, frontend-ul întâi cere proprietățile implicite de la
\texttt{ComponentDefaultsView}, apoi trimite un \texttt{POST} la
\texttt{ComponentUpsertView}. Orice schimbare de proprietate pleacă imediat la server,
tot printr-un \texttt{POST}, iar layoutul local se rescrie cu răspunsul. Serverul
rămâne sursa de adevăr, frontend-ul doar îl oglindește. Previzualizarea o randez
într-un iframe care afișează HTML-ul scos de motorul de export, ca să vezi
rezultatul real, așa cum va apărea după publicare.

Ca să coordoneze toate astea, editorul ține câteva variabile globale de stare.
\texttt{editorProject} și \texttt{editorLayout} sunt referințele curente la proiect și
la JSON-ul layoutului, resincronizate la fiecare răspuns al API-ului. \texttt{selectedId}
e componenta selectată acum: ea decide ce proprietăți apar în panoul din dreapta.
\texttt{isDirty} spune dacă ai modificări locale nesalvate; când e ridicat, pornește
\texttt{autoSaveTimer}, care trimite totul la server după câteva secunde de liniște,
fără să apeși tu „salvează". Și apoi sunt \texttt{undoStack} cu \texttt{redoStack},
istoricul local pe \texttt{Ctrl+Z} / \texttt{Ctrl+Y}. Aici e o dublare conștientă:
versionarea de pe server e plasa de siguranță durabilă, dar e prea lentă pentru un
„ups, anulează": nimeni nu vrea un request la server ca să dea undo. Stivele locale
fac undo/redo-ul instant, iar snapshot-urile rămân pentru istoricul serios.

Panoul din stânga ține cele 38 de tipuri în cinci categorii pliabile: Layout,
Conținut, Media, Interactive, Secțiuni. Click pe un tip, iar editorul cere
proprietățile implicite de la \texttt{ComponentDefaultsView} (un \texttt{hero}, de
pildă, vine gata cu titlu, subtitlu și două butoane, ca să nu te uiți la o cutie
goală), trimite upsert-ul și actualizează \texttt{editorLayout} cu răspunsul. Sus, în
bara de navigare, am pus trei butoane de previzualizare: desktop (100\%), tabletă
(768\,px), mobil (375\,px), care nu fac decât să redimensioneze iframe-ul, fără
reîncărcare. Așa verifici cum arată exportul pe telefon fără să ieși din editor și
fără să scoți de fapt arhiva.

\section{Considerații de securitate}

O aplicație cu autentificare și plăți atrage o anumită clasă de probleme, și am
încercat să le iau pe rând, nu să bifez o listă. Cel mai des întâlnit pericol e accesul
la resursa altcuiva, așa că prima linie de apărare e permisiunea \texttt{IsOwner} din
\texttt{core/permissions.py}: la fiecare cerere verifică \texttt{request.user ==
obj.owner} și taie accesul inter-utilizatori la proiecte, versiuni, pagini și submisii.
Important e că o aplic \emph{explicit} în fiecare ViewSet și APIView care atinge
resurse personale, nu mă bazez doar pe filtrarea implicită a query set-ului, fiindcă
aceea protejează listarea, dar nu și un acces direct pe ID.

A doua linie e validarea la frontiera API. \texttt{validate\_layout()} rulează înainte
de orice scriere pe \texttt{layout}, iar tipurile sunt închise la mulțimea
\texttt{VALID\_COMPONENT\_TYPES}. Asta acoperă două lucruri deodată: integritatea
datelor (nu intră structuri stricate) și tentativa de a strecura un tip inventat,
care altfel ar putea ajunge text arbitrar în HTML-ul exportat. E principiul de
validare la graniță pe care îl recomandă OWASP \cite{owasp2021}: nu ai încredere în
nimic din ce vine din afară.

Singurul endpoint deschis spre lume e cel de submisii de formulare, deci tot el e
singurul care are nevoie cu adevărat de rate limiting. Îl fac cu cache-ul Django: cheia se
compune din token-ul de partajare, ID-ul formularului și IP-ul clientului, adică un
contor separat pe fiecare formular și fiecare vizitator. Crește la fiecare cerere și,
peste zece pe oră, refuz. IP-ul îl iau din \texttt{X-Forwarded-For} dacă există (cazul
unui proxy sau load balancer în față), cu cădere pe \texttt{REMOTE\_ADDR}. Recunosc că
nu e blindat: cine schimbă IP-ul ușor trece de el, dar pentru spam-ul obișnuit de
boți e suficient, și nu am vrut să trag un Redis în proiect doar pentru atât.

La plăți, garda de idempotență oprește activarea dublă. Înainte să umblu la planul
utilizatorului, verific unicitatea lui \path{stripe_payment_intent_id} în tabela
\texttt{Payment}: dacă plata e deja înregistrată, răspund cu eroare și nu reexecut
nimic. Scenariul concret de care mă apăr e banal: omul dă refresh pe pagina de
confirmare, dar fără gardă l-ar fi taxat de două ori, ceea ce e genul de bug
pe care nu vrei să-l ai într-un flux de bani.

Restul sunt măsuri de configurație, dar nu mai puțin importante. CORS lasă să treacă
doar originile din \texttt{CORS\_ALLOWED\_ORIGINS}, cu
\texttt{CORS\_ALLOW\_CREDENTIALS=True} pentru header-ele de autorizare.
\texttt{SECRET\_KEY} e obligatoriu prin variabilă de mediu, fără valoare implicită:
prefer o aplicație care refuză să pornească decât una care pornește cu o cheie
default cunoscută. Iar \texttt{DEBUG} e \texttt{False} implicit, ca să nu scap din
greșeală pagini de eroare Django pline de detalii în producție.

% ═══════════════════════════════════════════════════════════════════════════
%  CAPITOLUL 4 - REZULTATE ȘI EVALUARE
% ═══════════════════════════════════════════════════════════════════════════
\chapter{Rezultate și Evaluare}

\section{Strategia de testare}

Am ales pentru IWB doar teste de integrare la nivel de API, fără să mock-uiesc baza de
date. E o decizie pe care merită să o explic, fiindcă nu e cea mai la modă. Testele
folosesc \texttt{APITestCase} din DRF, care ridică o bază de date de test izolată
pentru fiecare clasă, iar un client HTTP simulează cererile fără să deschidă un socket
real, deci suita întreagă rulează în secunde. Motivul pentru care am evitat
mock-urile pe baza de date e simplu: ele ascund tocmai bug-urile care mă speriau cel mai
tare, cele de migrare și de interacțiune cu ORM-ul, care apar numai pe o bază de date
adevărată. Un test care trece pe un mock și pică pe Postgres e mai rău decât niciun
test.

Peste tot revine același helper de autentificare: fiecare clasă își face un user cu
\texttt{User.objects.create\_user()}, scoate token-uri JWT printr-un \texttt{POST} la
\texttt{/api/auth/login/}, apoi cheamă
\texttt{self.client.credentials(HTTP\_AUTHORIZATION=f'Bearer \{access\}')} ca toate
cererile de după să fie autentificate. Cele mai importante mi se par testele de
izolare: creez doi utilizatori și verific că, atunci când unul cere resursa celuilalt,
primește \texttt{403 Forbidden} sau \texttt{404 Not Found}. Acolo se vede negru pe alb
dacă \texttt{IsOwner} chiar își face treaba.

\section{Acoperirea testelor}

În total sunt 68 de teste, împărțite inegal pe șase aplicații
(Tabelul~\ref{tab:tests}). Inegal, pentru că le-am pus acolo unde era riscul, nu ca
să iasă cifre rotunde pe fiecare modul.

\begin{table}[htbp]
  \centering
  \caption{Distribuția testelor pe aplicații}
  \label{tab:tests}
  \renewcommand{\arraystretch}{1.3}
  \begin{tabularx}{\textwidth}{>{\bfseries}l X r}
    \toprule
    Aplicație & Clase de teste acoperite & Nr. \\
    \midrule
    \texttt{landing}  & \texttt{LandingTests}, endpoint public, acces fără JWT & 2 \\
    \texttt{core}     & \texttt{ValidatorTests}, validare layout și componente & 11 \\
    \texttt{users}    & \texttt{RegisterTests}, \texttt{LoginTests},
                        \texttt{LogoutTests}, \texttt{ProfileTests} & 10 \\
    \texttt{projects} & \texttt{ProjectCRUDTests}, \texttt{ProjectVersionTests},
                        \texttt{FormSubmitTests}, \texttt{TagTests},
                        \texttt{PageTests}, \texttt{ProjectTemplateTests} & 29 \\
    \texttt{editor}   & \texttt{EditorLayoutTests}, layout, upsert, delete, reorder & 8 \\
    \texttt{exporter} & \texttt{ExportTests}, ZIP, JSON, auth, izolare utilizatori & 8 \\
    \midrule
    \textbf{Total} & & \textbf{68} \\
    \bottomrule
  \end{tabularx}
\end{table}

Cele 11 teste de pe \texttt{core} țintesc validarea, fiindcă acolo e zidul. Unul
singur e pozitiv: confirmă că un layout corect trece curat. Restul caută necazul:
\texttt{"id"} duplicat pe același nivel, \texttt{"id"} duplicat între părinte și copil
(cazul cross-nivel, cel pe care e cel mai ușor să-l ratezi), componentă fără
\texttt{"type"}, tip inventat (\texttt{"invalid\_type"}), un \texttt{"section"} fără
\texttt{"children"}, ID gol (\texttt{""}) și layout fără cheia rădăcină
\texttt{"components"}. Am vrut să ating fiecare ramură de \texttt{if} din
\texttt{validate\_component()} și \texttt{validate\_layout()}. Nu mă interesa că
validarea merge pe cazul fericit, ci că respinge fiecare fel de gunoi în parte.

\texttt{projects} adună cele mai multe teste, 29, pe șase clase. \texttt{ProjectCRUDTests}
acoperă create/list/update/delete și verifică faptul de bază: vezi doar proiectele
tale. \texttt{ProjectVersionTests} confirmă că fiecare operație de editor lasă un
snapshot și că restaurarea funcționează corect: după ea, layoutul e cel al versiunii
alese, iar starea de dinainte de restaurare e ea însăși salvată ca versiune nouă (fix
comportamentul „nu pierzi nimic" pe care l-am promis în Capitolul~3).
\texttt{FormSubmitTests} prinde și submisia normală, și rate limiting-ul: după zece
submisii reușite de pe același IP, a unsprezecea ia \texttt{429}. \texttt{TagTests},
\texttt{PageTests} și \texttt{ProjectTemplateTests} fac CRUD pe entitățile din jurul
proiectului.

\texttt{exporter} verifică partea care contează pentru utilizator: că arhiva ZIP are
chiar cele trei fișiere (\texttt{index.html}, \texttt{styles.css}, \texttt{main.js}),
că JSON-ul de previzualizare iese, că un neautentificat ia \texttt{401} și că B nu
poate exporta proiectul lui A. Toate cele 68 trec curat pe configurația de development,
pe SQLite. Nu pretind că 68 de teste înseamnă acoperire completă: ating endpoint-urile
și ramurile critice, dar mai sunt colțuri neacoperite, mai ales pe frontend, care nu e
testat automat deloc.

\section{Studii de caz}

\subsection{Cazul 1: Creare și export pagină simplă}

Un utilizator logat face un proiect minimal, \texttt{navbar} + \texttt{hero} +
\texttt{footer}, și exportă ZIP-ul. Cele trei fișiere (\texttt{index.html},
\texttt{styles.css}, \texttt{main.js}) merg fără nicio dependență: le deschizi direct
cu \texttt{file://} sau le pui pe orice server static și pagina arată la fel. Am
verificat asta inclusiv pe telefon: navbar-ul are hamburger funcțional, deși nu e
nici urmă de Bootstrap sau de vreun framework JS în arhivă. E cazul cel mai banal, dar
și cel care demonstrează promisiunea de bază: ce iese din IWB nu mai are nevoie de IWB.

\subsection{Cazul 2: Restaurarea unei versiuni}

Utilizatorul modifică layoutul de câteva ori și apoi vrea înapoi la o stare de mai
devreme, scenariul clasic „am stricat ceva, dă-mi varianta de acum o oră". Cheamă
endpoint-ul de restaurare cu ID-ul versiunii dorite; sistemul îngheață mai întâi starea
curentă într-un snapshot nou (cu o notă pusă automat) și abia apoi rescrie
\texttt{layout} cu cel ales. Subtilitatea, pe care o repet pentru că e ușor de ratat în
proiectele cu undo: nici starea curentă nu se pierde, ci devine și ea o versiune în
istoric. Așa poți merge înainte și înapoi printre toate stările, fără capăt de drum.

\subsection{Cazul 3: Partajare publică și colectare formulare}

Utilizatorul aprinde accesul public. De acolo, oricine are link-ul deschide pagina
fără cont, iar ce trimit vizitatorii prin formulare ajunge la proprietar prin
endpoint-ul de submisii. Spam-ul e ținut în frâu de rate limiting: a unsprezecea
submisie de pe același IP într-o oră ia \texttt{429 Too Many Requests}. Și dacă, dintr-un
motiv sau altul, vrei să tai accesul cuiva, regenerezi token-ul de partajare: vechiul
link moare instant, fără să fie nevoie să schimbi sau să ștergi proiectul.

\subsection{Cazul 4: Checkout Stripe și activare abonament}

Utilizatorul cere planul Pro. Frontend-ul ia \texttt{client\_secret} de la endpoint-ul
de creare a intenției și confirmă plata cu Stripe.js, direct pe card. La confirmarea de
server, întreb Stripe dacă \texttt{PaymentIntent}-ul e pe \texttt{succeeded}, trec prin
garda de idempotență și activez planul pe loc. Pleacă apoi emailul HTML cu detaliile:
sumă, referință, ultimele patru cifre ale cardului, data expirării. L-am testat repetat
cu cardurile de test Stripe, inclusiv dând refresh pe pagina de confirmare ca să mă
asigur că garda chiar oprește dubla activare; de fiecare dată planul s-a activat o
singură dată.

\section{Limitări ale sistemului}

Prefer să fiu sincer cu ce nu face IWB, fiindcă lista asta spune mai mult despre
proiect decât lauda. E un prototip de licență, și se vede.

Cea mai vizibilă limită e editorul în sine: nu am un canvas WYSIWYG cu drag-and-drop
adevărat. Adaugi și aranjezi componentele din panoul lateral, ceea ce funcționează, dar
e clar mai puțin intuitiv decât să tragi un bloc exact unde vrei pe pagină, cum face
Wix. Pentru o demonstrație de concept e suficient; pentru un utilizator real, ar fi
prima frustrare. Tot legat de editor: nu există editare simultană: doi oameni pe
același proiect și-ar călca în picioare, fiindcă nu am nicio strategie de rezolvare a
conflictelor.

Apoi sunt limitele de infrastructură. Imaginile uploadate stau pe disc local, ceea ce
e în regulă pe laptopul meu, dar cade din prima la un deployment serios: ai nevoie
de un S3 sau ceva similar. Rate limiting-ul îl am doar pe formulare, prin cache-ul
Django; în producție aș muta asta pe nginx sau Redis. Iar exportul, deși modelul de
date știe deja de pagini multiple, scoate deocamdată un singur \texttt{index.html}, fără
navigare între pagini; e o extindere pe care o am pe listă, nu o problemă de
arhitectură. Cele 38 de tipuri de componente sunt și ele puține față de sutele dintr-un
produs comercial, deși aici chiar e ușor de crescut: trei fișiere per tip nou, cum am
arătat.

% ═══════════════════════════════════════════════════════════════════════════
%  CAPITOLUL 5 - CONCLUZII
% ═══════════════════════════════════════════════════════════════════════════
\chapter{Concluzii și direcții viitoare de dezvoltare}

\section{Concluzii}

Lucrarea a urmărit, de la cap la coadă, cum se construiește IWB: un editor
web vizual de tip page builder, cu arhitectură client-server, backend Django~5.2 REST
și frontend SPA în JavaScript nativ.

Câteva decizii s-au dovedit centrale, și le-aș lua la fel a doua oară. Împărțirea
backend-ului în șapte aplicații cu granițe clare a făcut codul ușor de extins pe bucăți,
fără ca o modificare să rupă altceva. Validarea recursivă a layoutului la frontiera API
s-a dovedit zidul care chiar a ținut afară structurile stricate și tipurile inventate.
Versionarea automată, cu un snapshot înaintea fiecărei modificări, a transformat „am
pierdut munca" într-un non-eveniment. Iar motorul de export, partea de care sunt cel
mai mulțumit, livrează ce promite: un pachet HTML/CSS/JS care trăiește fără IWB,
pe orice hosting static, lucrul pe care platformele comerciale închise nu ți-l
dau. Peste astea, JWT-ul cu rotație și blacklisting taie reutilizarea token-urilor
furate, iar integrarea Stripe încasează bani reali cu o gardă de idempotență care nu
lasă să se întâmple dubla activare.

Dacă e o concluzie pe care vreau să o rețină cineva, e că nu ai nevoie de un stack
greu ca să faci asta. Un page builder funcțional încape pe tehnologii open-source
standard, fără React și fără infrastructură scumpă, iar cele 68 de teste de
integrare, rulate pe o bază de date reală, arată că lanțul ține în condiții
apropiate de cele de producție.

\section{Direcții viitoare de dezvoltare}

Dacă ar fi să continui proiectul, și sper să o fac, prioritățile nu sunt egale.
Le-aș pune cam în ordinea de mai jos, de la ce schimbă cel mai mult experiența până la
ce întărește colțurile.

\subsection{Canvas drag-and-drop vizual}
Asta e, fără discuție, lucrul care lipsește cel mai tare. Un canvas pe care tragi
blocul exact unde vrei, în timp real, ar apropia IWB de senzația unui editor comercial.
Nu e trivial: ar cere un management de stare mai serios în frontend și o sincronizare
mai fină cu backend-ul decât „trimite tot layoutul la fiecare schimbare". Probabil aici
ar ceda în sfârșit decizia mea de a evita un framework.

\subsection{Export multi-pagină}
Următoarea pe listă, și cea mai aproape de a fi gata. Modelul de date știe deja de
pagini multiple (\texttt{Page} are câmp \texttt{order}), deci nu mai rămâne decât
ca motorul de export să itereze toate paginile și să scoată câte un HTML pentru fiecare,
cu navigare între ele. E mai mult muncă de scris decât problemă de gândit.

\subsection{Colaborare în timp real}
Editare simultană, mai mulți oameni pe același proiect. Ar însemna WebSocket (Django
Channels) și o strategie reală de rezolvare a conflictelor, pe CRDT sau transformări
operaționale \cite{shapiro2011}. Partea grea, de fapt, e a doua, nu transportul.

\subsection{Stocarea imaginilor în cloud}
Mutarea imaginilor de pe discul local pe S3, Cloudflare R2 sau echivalent. E condiția
fără de care un deployment scalabil nici nu are sens, și e și relativ ușor de făcut:
Django are deja backend-uri de storage pentru asta.

\subsection{Teme și personalizare vizuală}
Un sistem de teme (palete, fonturi) care să atingă tot site-ul exportat printr-un
singur set de variabile. Aici infrastructura e deja pusă: \texttt{meta} poate ține
variabilele CSS globale, iar motorul de export le injectează. Mai lipsește doar
interfața care să le aleagă frumos.

\subsection{Indexare SEO și webhooks Stripe}
Două colțuri de întărit. Pe de o parte, un pre-render sau SSR pentru ca paginile
publice să fie indexabile, ceea ce rezolvă dezavantajul clasic al unui SPA. Pe de alta, un endpoint
de webhook Stripe pentru evenimentele asincrone (eșec de plată, anulare de abonament),
ca fluxul de bani să nu mai depindă doar de confirmarea care vine de la client și să
reziste la un timeout de rețea.

% ═══════════════════════════════════════════════════════════════════════════
%  BIBLIOGRAFIE (inline - nu necesită fișier .bib separat)
% ═══════════════════════════════════════════════════════════════════════════
\begin{thebibliography}{99}
\addcontentsline{toc}{chapter}{Bibliografie}

\bibitem{djangodocs}
Django Software Foundation.
\textit{Django Documentation, Version 5.2}, 2025.
\url{https://docs.djangoproject.com/en/5.2/}.
Accesat: mai 2026.

\bibitem{drfdocs}
T. Christie and contributors.
\textit{Django REST Framework Documentation}, 2024.
\url{https://www.django-rest-framework.org/}.
Accesat: mai 2026.

\bibitem{simplejwt}
Jazzband contributors.
\textit{Simple JWT, A JSON Web Token authentication plugin for Django REST Framework}, 2024.
\url{https://django-rest-framework-simplejwt.readthedocs.io/}.
Accesat: mai 2026.

\bibitem{stripedocs}
Stripe, Inc.
\textit{Stripe API Reference, Python SDK}, 2025.
\url{https://stripe.com/docs/api}.
Accesat: mai 2026.

\bibitem{sqlitedocs}
R. Hipp and contributors.
\textit{SQLite Documentation}, 2024.
\url{https://www.sqlite.org/docs.html}.
Accesat: mai 2026.

\bibitem{postgresqldocs}
PostgreSQL Global Development Group.
\textit{PostgreSQL 16 Documentation}, 2024.
\url{https://www.postgresql.org/docs/16/}.
Accesat: mai 2026.

\bibitem{rfc7519}
M. Jones, J. Bradley, and N. Sakimura.
\textit{RFC 7519: JSON Web Token (JWT)}.
Internet Engineering Task Force (IETF), May 2015.
\url{https://www.rfc-editor.org/rfc/rfc7519}.

\bibitem{rfc8259}
T. Bray.
\textit{RFC 8259: The JavaScript Object Notation (JSON) Data Interchange Format}.
IETF, December 2017.
\url{https://www.rfc-editor.org/rfc/rfc8259}.

\bibitem{fielding2000}
R. T. Fielding.
\textit{Architectural Styles and the Design of Network-based Software Architectures}.
PhD thesis, University of California, Irvine, 2000.
\url{https://ics.uci.edu/~fielding/pubs/dissertation/top.htm}.

\bibitem{shapiro2011}
M. Shapiro, N. Preguiça, C. Baquero, and M. Zawirski.
Conflict-Free Replicated Data Types.
In \textit{Stabilization, Safety, and Security of Distributed Systems (SSS 2011)},
pages 386--400. Springer, 2011.

\bibitem{mdnSPA}
Mozilla Developer Network contributors.
\textit{Single-page application (SPA), MDN Web Docs Glossary}, 2024.
\url{https://developer.mozilla.org/en-US/docs/Glossary/SPA}.
Accesat: mai 2026.

\bibitem{mdnFetch}
Mozilla Developer Network contributors.
\textit{Fetch API, Web APIs, MDN Web Docs}, 2024.
\url{https://developer.mozilla.org/en-US/docs/Web/API/Fetch_API}.
Accesat: mai 2026.

\bibitem{owasp2021}
OWASP Foundation.
\textit{OWASP Top Ten 2021}, 2021.
\url{https://owasp.org/www-project-top-ten/}.
Accesat: mai 2026.

\bibitem{wixwebsite}
Wix.com.
\textit{Wix, Create a Free Website}, 2024.
\url{https://www.wix.com/}.
Accesat: mai 2026.

\bibitem{webflowsite}
Webflow, Inc.
\textit{Webflow: Create a custom website}, 2024.
\url{https://webflow.com/}.
Accesat: mai 2026.

\bibitem{squarespace}
Squarespace, Inc.
\textit{Squarespace, Build a Website}, 2024.
\url{https://www.squarespace.com/}.
Accesat: mai 2026.

\bibitem{grapesjs}
A. Arseniev and contributors.
\textit{GrapesJS, Free and Open Source Web Builder Framework}, 2024.
\url{https://grapesjs.com/}.
Accesat: mai 2026.

\bibitem{w3techs}
W3Techs.
\textit{Usage Statistics and Market Share of Content Management Systems}, 2025.
\url{https://w3techs.com/technologies/overview/content_management}.
Accesat: mai 2026.

\end{thebibliography}

\end{document}
